\documentclass[11pt, a4paper]{article}
\usepackage{amsmath}
\usepackage{amssymb}
\usepackage{graphicx}
\usepackage{hyperref}
\usepackage{xcolor}
\usepackage{listings}
\usepackage[margin=1in]{geometry}
\usepackage{fancyhdr}
\usepackage{titlesec}
\usepackage{booktabs}
\usepackage{enumitem}
\usepackage{tabularx}

\lstset{
    basicstyle=\ttfamily\small,
    keywordstyle=\color{blue},
    commentstyle=\color{gray},
    stringstyle=\color{red},
    breaklines=true,
    language=Python
}

\title{ECG Simulator Documentation\\Part 4: GUI Reference \& Parameter Reference}
\author{Comprehensive Technical Documentation}
\date{June 2026}

\pagestyle{fancy}
\fancyhf{}
\rhead{Part 4: GUI Reference}
\lhead{ECG Simulator}
\cfoot{\thepage}

\begin{document}

\maketitle

\tableofcontents
\newpage

\section{Main Window Architecture}

\subsection{Overview}
The PyQt6-based GUI provides real-time visualization of 12-lead ECG with interactive parameter control. The main window is organized into three primary regions:

\begin{itemize}
    \item \textbf{Control Panel} (left): Heart rate, engine selection, beat controls
    \item \textbf{ECG Display} (center): Real-time scrolling 12-lead ECG trace
    \item \textbf{Parameter Panel} (right): 50+ physiological parameter controls
\end{itemize}

\subsection{Control Panel Layout}

\subsubsection{Heart Rate (HR) Control}
\begin{itemize}
    \item \textbf{Widget Type}: QSpinBox
    \item \textbf{Range}: 30--220 bpm
    \item \textbf{Default}: 70 bpm
    \item \textbf{Step}: 1 bpm
    \item \textbf{Effect}: Directly sets SA node cycle length via:
    \begin{equation}
        \text{cycle\_length} = \frac{60}{HR}
    \end{equation}
    \item \textbf{Clinical Range}: 60--100 bpm (normal sinus rhythm)
    \item \textbf{Bradycardia}: <60 bpm (typically 40--60 in athletes)
    \item \textbf{Tachycardia}: >100 bpm (up to 200+ during exercise)
\end{itemize}

\subsubsection{Simulation Engine Selector}
\begin{itemize}
    \item \textbf{Widget Type}: QComboBox
    \item \textbf{Options}: 
    \begin{itemize}
        \item \texttt{graph} (Phase 1/2: BFS-based, >1000 Hz, deterministic)
        \item \texttt{cable} (Phase 3: FHN cable, 2.2× real-time, physiologically accurate)
    \end{itemize}
    \item \textbf{Default}: \texttt{graph}
    \item \textbf{Effect}: Switches underlying simulation engine
    \item \textbf{Behavior}: Resets current beat on switch
\end{itemize}

\subsubsection{Beat Control Buttons}
\begin{itemize}
    \item \textbf{Start Button}: Begin continuous rhythm generation
    \begin{itemize}
        \item Starts simulation timer (33 ms per frame @ 30 FPS)
        \item Disables beat controls during active simulation
    \end{itemize}
    \item \textbf{Stop Button}: Pause simulation (preserves current ECG trace)
    \begin{itemize}
        \item Allows parameter adjustment without clearing ECG
    \end{itemize}
    \item \textbf{Reset Button}: Clear ECG trace and restart
    \begin{itemize}
        \item Clears all accumulated beats from display
        \item Resets internal beat counters
    \end{itemize}
\end{itemize}

\section{ECG Display Widget}

\subsection{Real-Time 12-Lead Display}
\begin{itemize}
    \item \textbf{Library}: pyqtgraph (for real-time plotting performance)
    \item \textbf{Update Rate}: 30 FPS
    \item \textbf{Scroll Behavior}: Left-to-right, continuous scroll
    \item \textbf{Time Window}: ~10 seconds visible (adjustable)
    \item \textbf{Leads Displayed}: All 12 ECG leads (6 frontal + 6 precordial)
\end{itemize}

\subsection{12-Lead Organization}

\subsubsection{Frontal Leads (rows 1--2)}
\begin{tabularx}{\textwidth}{|X|X|X|}
\hline
\textbf{Lead} & \textbf{Axis} & \textbf{Clinical Use} \\
\hline
I & 0° (left) & Left-right forces \\
II & 60° (inferolateral) & Inferior/lateral forces \\
III & 120° (inferior) & Inferior forces \\
aVR & -150° (right) & Right atrium/ventricle \\
aVL & -30° (left) & Left ventricle \\
aVF & 90° (inferior) & Inferior wall \\
\hline
\end{tabularx}

\subsubsection{Precordial Leads (rows 3--4)}
\begin{tabularx}{\textwidth}{|X|X|X|}
\hline
\textbf{Lead} & \textbf{Anatomical Position} & \textbf{Chamber} \\
\hline
V1 & 4th ICS, right sternal border & RV outflow \\
V2 & 4th ICS, left sternal border & LV outflow \\
V3 & Midway V2--V4 & LV septum \\
V4 & 5th ICS, midclavicular & LV apex \\
V5 & 5th ICS, anterior axillary & LV lateral \\
V6 & 5th ICS, midaxillary & LV lateral \\
\hline
\end{tabularx}

\section{Parameter Panel Organization}

\subsection{Structure}
The parameter panel is organized into 12 collapsible groups:

\begin{enumerate}
    \item SA Node Parameters
    \item AV Node Parameters
    \item His Bundle \& Purkinje
    \item Atrial Conduction
    \item Ventricular Conduction
    \item Ventricular Refractory Periods
    \item Escape Pacemaker
    \item Atrial Fibrillation
    \item Ectopic Focus
    \item ECG Scaling
    \item Plugin Controls
    \item Advanced/Debug
\end{enumerate}

\newpage
\section{Complete Parameter Reference (50+ Parameters)}

\subsection{Group 1: SA Node Parameters}

\begin{tabularx}{\textwidth}{|l|X|c|c|}
\hline
\textbf{Parameter} & \textbf{Description} & \textbf{Range} & \textbf{Unit} \\
\hline
cycle\_length & Interval between SA fires & 0.3--2.0 & s \\
funny\_current & HCN4 channel conductance & 0.0--1.0 & g/g\_max \\
autonomic\_tone & Sympathetic/parasympathetic balance & -1.0 to +1.0 & normalized \\
\hline
\end{tabularx}

\subsubsection{Usage Notes}
\begin{itemize}
    \item \textbf{cycle\_length}: Typically set via HR spinner; direct control available
    \item \textbf{funny\_current}: Alters SA pacemaker rate
        \begin{itemize}
            \item 0.0 = No spontaneous depolarization (requires external stimulus)
            \item 1.0 = Maximum rate (~200 bpm)
        \end{itemize}
    \item \textbf{autonomic\_tone}: Modulates HR via sympathetic/parasympathetic effects
        \begin{itemize}
            \item -1.0 = Maximum parasympathetic (vagal): ~40 bpm
            \item 0.0 = Balanced (normal): ~70 bpm
            \item +1.0 = Maximum sympathetic: ~120 bpm
        \end{itemize}
\end{itemize}

\subsection{Group 2: AV Node Parameters}

\begin{tabularx}{\textwidth}{|l|X|c|c|}
\hline
\textbf{Parameter} & \textbf{Description} & \textbf{Range} & \textbf{Unit} \\
\hline
av\_delay & Conduction time through AV node & 0.05--0.30 & s \\
av\_refractory & AV node absolute refractory period & 0.20--0.40 & s \\
conductance & AV nodal conductance (0=block) & 0.0--1.0 & g/g\_max \\
\hline
\end{tabularx}

\subsubsection{Clinical Significance}
\begin{itemize}
    \item \textbf{av\_delay}: Normal = 100--200 ms
        \begin{itemize}
            \item Controls PR interval via: $PR = P\text{-wave} + av\_delay + QRS$
            \item Prolonged PR (First-degree AV block) when >200 ms
        \end{itemize}
    \item \textbf{conductance}: Controls AV block severity
        \begin{itemize}
            \item 1.0 = Normal conduction (1:1 conduction ratio)
            \item 0.5 = 2:1 AV block (every 2nd atrial beat blocked)
            \item 0.0 = Complete AV block (III° AV block)
        \end{itemize}
\end{itemize}

\subsection{Group 3: His Bundle \& Purkinje System}

\begin{tabularx}{\textwidth}{|l|X|c|c|}
\hline
\textbf{Parameter} & \textbf{Description} & \textbf{Range} & \textbf{Unit} \\
\hline
his\_delay & His bundle conduction time & 0.01--0.05 & s \\
lbb\_conductance & Left bundle branch conductance & 0.0--1.0 & g/g\_max \\
rbb\_conductance & Right bundle branch conductance & 0.0--1.0 & g/g\_max \\
lahb\_conductance & Left ant. fascicle conductance & 0.0--1.0 & g/g\_max \\
lphb\_conductance & Left post. fascicle conductance & 0.0--1.0 & g/g\_max \\
\hline
\end{tabularx}

\subsubsection{Bundle Branch Block Patterns}
\begin{itemize}
    \item \textbf{LBBB} (Left Bundle Branch Block):
    \begin{itemize}
        \item Set lbb\_conductance = 0.0, keep rbb\_conductance = 1.0
        \item ECG: Broad QRS (>120 ms), M-shaped in V5--V6
    \end{itemize}
    \item \textbf{RBBB} (Right Bundle Branch Block):
    \begin{itemize}
        \item Set rbb\_conductance = 0.0, keep lbb\_conductance = 1.0
        \item ECG: Broad QRS (>120 ms), M-shaped in V1--V2
    \end{itemize}
    \item \textbf{LAHB} (Left Anterior Hemiblock):
    \begin{itemize}
        \item Set lahb\_conductance = 0.0, keep lphb\_conductance = 1.0
        \item ECG: Left axis deviation, QRS <120 ms
    \end{itemize}
    \item \textbf{LPHB} (Left Posterior Hemiblock):
    \begin{itemize}
        \item Set lphb\_conductance = 0.0, keep lahb\_conductance = 1.0
        \item ECG: Right axis deviation, QRS <120 ms
    \end{itemize}
\end{itemize}

\subsection{Group 4: Atrial Conduction}

\begin{tabularx}{\textwidth}{|l|X|c|c|}
\hline
\textbf{Parameter} & \textbf{Description} & \textbf{Range} & \textbf{Unit} \\
\hline
ra\_conduction & Right atrial conduction velocity & 0.50--1.50 & m/s \\
la\_conduction & Left atrial conduction velocity & 0.50--1.50 & m/s \\
atrial\_refractory & Atrial myocardium absolute refractory & 0.20--0.35 & s \\
\hline
\end{tabularx}

\subsection{Group 5: Ventricular Conduction}

\begin{tabularx}{\textwidth}{|l|X|c|c|}
\hline
\textbf{Parameter} & \textbf{Description} & \textbf{Range} & \textbf{Unit} \\
\hline
sept\_conduction & Septal conduction velocity & 0.3--1.5 & m/s \\
lv\_conduction & LV free wall conduction velocity & 0.30--1.0 & m/s \\
rv\_conduction & RV free wall conduction velocity & 0.30--1.0 & m/s \\
\hline
\end{tabularx}

\subsubsection{QRS Duration Formula}
Total QRS duration depends on conduction velocities and geometric path lengths.
For Phase 1/2 engine:
\begin{equation}
    QRS_{\text{duration}} \approx \frac{\text{longest path}}{v_{\text{min}}}
\end{equation}

\subsection{Group 6: Ventricular Refractory Periods}

\begin{tabularx}{\textwidth}{|l|X|c|c|}
\hline
\textbf{Parameter} & \textbf{Description} & \textbf{Range} & \textbf{Unit} \\
\hline
vent\_abs\_refr & Absolute refractory period & 0.25--0.45 & s \\
vent\_rel\_refr & Relative refractory period & 0.10--0.35 & s \\
\hline
\end{tabularx}

\subsubsection{Physiological Basis}
\begin{itemize}
    \item \textbf{Absolute Refractory Period (ARP)}: Voltage-gated Na$^+$ channels inactivated
        \begin{itemize}
            \item No stimulus can trigger action potential
            \item Normal duration: 200--300 ms
            \item Set on Phase 1/2 engine only
        \end{itemize}
    \item \textbf{Relative Refractory Period (RRP)}: Partial recovery of Na$^+$ channel availability
        \begin{itemize}
            \item Stimulus can trigger AP but with reduced amplitude
            \item Conduction may be slower or blocked
            \item Normal duration: 150--300 ms additional
        \end{itemize}
    \item \textbf{Phase 3 Engine}: Refractory periods emerge naturally from FHN $w$ variable
        \begin{itemize}
            \item Recovery time determined by:
            \begin{equation}
                \tau_w = \frac{\tau_s}{b \cdot \varepsilon} = \frac{0.015}{0.8 \times 0.08} \approx 234 \text{ ms}
            \end{equation}
            \item No explicit refractory parameter needed
        \end{itemize}
\end{itemize}

\subsection{Group 7: Escape Pacemaker}

\begin{tabularx}{\textwidth}{|l|X|c|c|}
\hline
\textbf{Parameter} & \textbf{Description} & \textbf{Range} & \textbf{Unit} \\
\hline
escape\_enabled & Enable backup pacemaker & Yes/No & bool \\
escape\_interval & Escape interval (interval without SA) & 0.5--3.0 & s \\
escape\_origin & Ectopic site (AV/HIS/LV) & enum & - \\
\hline
\end{tabularx}

\subsubsection{Escape Pacemaker Mechanism}
\begin{itemize}
    \item Automatically fires if no atrial impulse received
    \item Clinical significance:
    \begin{itemize}
        \item \textbf{Junctional escape} (origin=AV): Fires at ~40--60 bpm, narrow QRS
        \item \textbf{Ventricular escape} (origin=LV): Fires at ~20--40 bpm, broad QRS
    \end{itemize}
    \item escape\_interval determines delay before firing:
    \begin{itemize}
        \item If no stimulus received for >escape\_interval seconds, escape fires
        \item Common clinical values: 1.5--3.0 seconds
    \end{itemize}
\end{itemize}

\subsection{Group 8: Atrial Fibrillation}

\begin{tabularx}{\textwidth}{|l|X|c|c|}
\hline
\textbf{Parameter} & \textbf{Description} & \textbf{Range} & \textbf{Unit} \\
\hline
af\_enabled & Enable atrial fibrillation & Yes/No & bool \\
af\_rate & F-wave (fibrillatory wave) rate & 300--700 & bpm \\
af\_rr\_distribution & RR interval variation type & fixed/poisson/gaussian & enum \\
af\_f\_amplitude & F-wave amplitude & 0.01--0.20 & mV \\
\hline
\end{tabularx}

\subsubsection{AF ECG Characteristics}
\begin{itemize}
    \item \textbf{AF Rate}: Normal 300--600 bpm (fibrillatory waves from disorganized reentry)
        \begin{itemize}
            \item Slower rates (300 bpm) suggest possible organized activity (flutter)
            \item Faster rates (600+ bpm) suggest rapid fibrillatory activity
        \end{itemize}
    \item \textbf{RR Distribution}:
    \begin{itemize}
        \item fixed: Constant interval (used for testing)
        \item poisson: Random exponential intervals (physiologic, irregular)
        \item gaussian: Normal distribution around mean (moderate irregularity)
    \end{itemize}
    \item \textbf{F-wave Amplitude}: Visible baseline oscillations
        \begin{itemize}
            \item Coarse AF: >0.10 mV (more organized)
            \item Fine AF: <0.05 mV (more disorganized)
        \end{itemize}
\end{itemize}

\subsection{Group 9: Ectopic Focus}

\begin{tabularx}{\textwidth}{|l|X|c|c|}
\hline
\textbf{Parameter} & \textbf{Description} & \textbf{Range} & \textbf{Unit} \\
\hline
ectopic\_enabled & Enable ectopic focus & Yes/No & bool \\
ectopic\_coupling & Coupling interval (PVC timing) & 0.40--0.80 & s \\
ectopic\_repetitive & Consecutive ectopics & Yes/No & bool \\
ectopic\_origin & Ectopic site & LA/RA/LV/RV & enum \\
\hline
\end{tabularx}

\subsubsection{Premature Contraction Types}
\begin{itemize}
    \item \textbf{PAC (Premature Atrial Contraction)}:
    \begin{itemize}
        \item Set ectopic\_origin = LA or RA
        \item Normal coupling: 350--450 ms
        \item Early coupling (<350 ms) may block AV node → aberrant conduction
    \end{itemize}
    \item \textbf{PVC (Premature Ventricular Contraction)}:
    \begin{itemize}
        \item Set ectopic\_origin = LV or RV
        \item Normal coupling: 400--600 ms
        \item ECG: Broad QRS (>120 ms), T-wave opposite normal polarity
        \item Couplets/triplets possible if ectopic\_repetitive = true
    \end{itemize}
\end{itemize}

\subsection{Group 10: ECG Scaling}

\begin{tabularx}{\textwidth}{|l|X|c|c|}
\hline
\textbf{Parameter} & \textbf{Description} & \textbf{Range} & \textbf{Unit} \\
\hline
amplitude\_scale & Global amplitude multiplier & 0.1--5.0 & unitless \\
baseline\_shift & DC offset for visibility & -2.0--+2.0 & mV \\
\hline
\end{tabularx}

\newpage
\section{Advanced Controls \& Debug Parameters}

\subsection{Group 11: Plugin Controls}
\begin{itemize}
    \item \textbf{Active Plugins}: Dropdown showing enabled pathologies
    \item \textbf{Plugin Stacking}: Select multiple pathologies simultaneously
    \item \textbf{Save/Load Presets}: Store/retrieve parameter combinations
\end{itemize}

\subsection{Group 12: Advanced/Debug}

\begin{tabularx}{\textwidth}{|l|X|c|c|}
\hline
\textbf{Parameter} & \textbf{Description} & \textbf{Range} & \textbf{Unit} \\
\hline
debug\_mode & Enable diagnostic output & Yes/No & bool \\
show\_activation\_times & Print node activation times & Yes/No & bool \\
cable\_internal\_dt & Phase 3 integration step & 0.0001--0.001 & s \\
\hline
\end{tabularx}

\section{Parameter Adjustment Workflow}

\subsection{Typical Use Cases}

\subsubsection{Scenario 1: Simulate Normal Sinus Rhythm}
\begin{enumerate}
    \item HR = 70 bpm
    \item All conductances = 1.0
    \item All refractories = defaults
    \item af\_enabled = false
    \item ectopic\_enabled = false
\end{enumerate}

\subsubsection{Scenario 2: Simulate First-Degree AV Block}
\begin{enumerate}
    \item HR = 70 bpm
    \item av\_delay = 0.25 s (prolonged from normal 0.15 s)
    \item av\_conductance = 1.0 (1:1 conduction maintained)
    \item Result: PR interval >200 ms
\end{enumerate}

\subsubsection{Scenario 3: Simulate Left Bundle Branch Block}
\begin{enumerate}
    \item HR = 70 bpm
    \item lbb\_conductance = 0.0 (complete block)
    \item rbb\_conductance = 1.0 (normal)
    \item Result: Broad QRS (>120 ms), wide S wave in V5--V6
\end{enumerate}

\subsubsection{Scenario 4: Simulate Atrial Fibrillation}
\begin{enumerate}
    \item HR = 130 bpm (controlled ventricular rate)
    \item af\_enabled = true
    \item af\_rate = 450 bpm (fibrillatory waves)
    \item af\_rr\_distribution = poisson
    \item Result: Irregular rhythm, absent P waves, f-waves visible
\end{enumerate}

\section{Real-Time Display Updates}

\subsection{Refresh Mechanism}
\begin{itemize}
    \item GUI updates at 30 FPS
    \item ECG trace accumulates every ~33 ms
    \item Leads redrawn continuously
    \item Parameters can be adjusted without stopping simulation
\end{itemize}

\subsection{Performance Notes}
\begin{itemize}
    \item Phase 1/2 engine: >1000 Hz update rate → GPU-limited by display
    \item Phase 3 engine: ~150 Hz (2.2× real-time) → smooth real-time performance
    \item Parameter changes take effect immediately on next beat
    \item HR changes cause immediate cycle\_length update
\end{itemize}

\end{document}
